\chapter{多重网格操作的并行结构}
\label{ch:mg-parallel-operations}

\begin{keybox}
\textbf{本章把\chapref{ch:parallel-model}的抽象模型具体化到多重网格。}对每一个 smoother、算子应用、残差、restriction、prolongation、correction、norm 与 ghost fill，都沿着同一条链分析：
\begin{equation}
\boxed{
\text{数值语义}
\rightarrow
\text{状态版本}
\rightarrow
\text{logical task}
\rightarrow
\text{依赖 DAG}
\rightarrow
\text{所有权与粒度}
\rightarrow
\text{同步/通信来源}.
}
\label{eq:mg-parallel-analysis-chain}
\end{equation}
本章仍不选择 OpenMP、CUDA 或 MPI。目标是先回答一个更基础的问题：\textbf{V-cycle 中究竟哪些依赖是数学上必须存在的，哪些依赖只是某种存储或执行方式引入的，以及粗化为什么会系统性地耗尽并行度。}
\end{keybox}

\section{一个 MG Level 上究竟有哪些状态}

对线性多重网格的第 $\ell$ 层，最基本的逻辑对象可以写成
\begin{equation}
\boxed{
u_\ell,\quad
b_\ell,\quad
r_\ell,\quad
e_\ell,\quad
A_\ell,\quad
\text{geometry/coefficient data}.
}
\label{eq:mg-parallel-level-data}
\end{equation}
其中 $u_\ell$ 是当前近似，$b_\ell$ 是当前层右端，
\begin{equation}
r_\ell=b_\ell-A_\ell u_\ell
\end{equation}
是残差，$e_\ell$ 表示误差修正。若采用无矩阵算子，$A_\ell$ 不必以 CSR 等显式稀疏矩阵形式存储，但应用 $A_\ell$ 所需的 stencil、面系数、几何信息与边界元数据仍必须存在。

并行分析时，仅列出这些数组还不够。真正决定依赖的是\textbf{状态版本}。例如一次 Jacobi sweep 中，
\begin{equation}
u_\ell^{(k)}
\longrightarrow
u_\ell^{(k+1)}
\end{equation}
是两个不同的逻辑状态；即使实现时通过 ping-pong buffer 复用物理内存，也不能把二者在依赖图中混成同一个对象。相反，Gauss--Seidel 会在同一个 sweep 内逐步产生并立即消费新状态，因此其 DAG 与 Jacobi 根本不同。

这给出本章第一个重要区分：
\begin{equation}
\boxed{
\text{逻辑状态版本}
\neq
\text{物理数组对象}.
}
\label{eq:mg-state-version-vs-storage}
\end{equation}
前者决定算法依赖，后者决定内存占用和可能出现的 WAR/WAW。许多“并行化困难”其实来自把逻辑上不同的版本强行复用到同一块存储。

在结构网格中，状态通常又以 patch、box 或连续子块为存储与调度单位。热点索引可能采用某一方向连续的一维布局，例如三维 $x$ 连续时
\begin{equation}
p(i,j,k)=i+n_x\bigl[(j-1)+n_y(k-1)\bigr].
\end{equation}
布局本身不改变数学 DAG，却决定 task 聚合以后是否具有良好的 cache line、SIMD 或 coalesced memory access。

\section{把第七章模型落到 MG 操作}

对后面的每一个操作，不再重新发明一套问题，而是直接使用\chapref{ch:parallel-model}的分析流程。最核心的六步可以压缩为：

\begin{enumerate}
\item 写清楚这一操作要产生的\textbf{新状态}；
\item 选择最小有意义的 logical task，并写出 $R_t,W_t$；
\item 判断 RAW 真依赖在哪里，哪些 WAR/WAW 只是原地存储引入；
\item 判断 task 能否重写成 unique-writer 的 gather 形式；
\item 决定实际调度粒度是 point、line、tile、patch 还是更大的 block；
\item 找出必须等待的数据 readiness：本地阶段完成、ghost 就绪、coarse state 就绪或 global reduction 完成。
\end{enumerate}

这样分析以后，OpenMP barrier、GPU kernel boundary、MPI halo 或 collective 都只是同一条依赖边在不同机器上的物理实现，而不是彼此无关的“特殊技巧”。

\section{平滑器首先改变依赖图}

\subsection{Jacobi 的并行性来自状态分离}

Weighted Jacobi 的完整更新是
\begin{equation}
u_i^{(k+1)}
=
u_i^{(k)}
+
\omega D_{ii}^{-1}
\left(
b_i-(Au^{(k)})_i
\right).
\label{eq:mg-parallel-jacobi-full}
\end{equation}
一个 logical task 负责一个新输出 $u_i^{(k+1)}$。它只读取旧状态 $u^{(k)}$ 的 stencil 邻域、$b_i$ 与系数，写入自己的新状态：
\begin{equation}
R_i=
\{u_j^{(k)}:j\in\mathcal N(i)\}
\cup
\{b_i,\text{coefficients}\},
\qquad
W_i=\{u_i^{(k+1)}\}.
\label{eq:mg-jacobi-rw}
\end{equation}
不同 task 可以共享旧输入，却没有输出写冲突，因此同一 sweep 的 point tasks 在算法层面完全独立。

这里有两种实现语义必须分清。若先显式计算
\begin{equation}
r^{(k)}=b-Au^{(k)},
\end{equation}
再做
\begin{equation}
u^{(k+1)}=u^{(k)}+\omega D^{-1}r^{(k)},
\end{equation}
则 residual 是 stencil gather，update 是 map，中间存在 $r^{(k)}$ readiness。若把二者融合，单个 task 直接按式~\eqref{eq:mg-parallel-jacobi-full}产生 $u_i^{(k+1)}$，可以省去一遍 residual 数组的写入与读取，但会改变数据生命周期和 kernel 粒度。

Jacobi 容易并行的根本原因不是“公式简单”，而是
\begin{equation}
\boxed{
u^{(k)}\ \text{在整个 sweep 内只读},
\qquad
u^{(k+1)}\ \text{唯一写入}.
}
\end{equation}
double buffering 把原本可能由原地复用产生的 WAR/WAW 去掉了。

\subsection{Gauss--Seidel 的并行障碍来自新值立即可见}

Lexicographic Gauss--Seidel 在更新 $u_i$ 时会读取本 sweep 已经产生的某些 $u_j^{(k+1)}$。因此其依赖不是存储偶然造成的，而是更新规则本身要求“新值立即进入后续计算”。在最直接排序下，DAG 会出现沿网格传播的长链。

对于标准二维五点或三维七点 Cartesian stencil，耦合图是二分图，可以采用 red--black Gauss--Seidel。一个 sweep 变成
\begin{equation}
\boxed{
\text{all red}
\rightarrow
\text{red state ready}
\rightarrow
\text{all black}
\rightarrow
\text{black state ready}.
}
\label{eq:mg-rbgs-dag}
\end{equation}
同色内部 task 可并行，不同颜色之间存在 RAW 真依赖。

需要保留这个条件的边界：\textbf{“两色即可”来自 stencil 耦合图的二分性，而不是 Gauss--Seidel 的普遍性质。}若离散产生对角耦合、非结构邻接或 block coupling，可能需要更多颜色，或者改用 block/line smoother。颜色数增加会同时增加并行阶段数和同步次数。

\subsection{Line 与 Block Smoother 改变的是 task 粒度}

各向异性问题常使用 line smoother 或 block smoother。此时一个 logical task 不再是单点，而可能是“一条线上的小线性系统”或“一个 block 内的局部求解”。

这类方法常以更多 task 内部 work 换取更强的平滑效果：
\begin{equation}
\text{larger local solve}
\quad\Longleftrightarrow\quad
\text{fewer, heavier tasks}.
\end{equation}
因此不能只比较“每 sweep 的并行度”。真正需要比较的是
\begin{equation}
\boxed{
\text{time per sweep}
\times
\text{sweeps/cycles required for target convergence}.
}
\end{equation}
算法收敛性与并行粒度在 smoother 上第一次直接耦合起来。

\section{算子应用与残差是典型 Stencil Gather}

残差定义为
\begin{equation}
r_i=b_i-(Au)_i.
\end{equation}
对显式 stencil 或无矩阵 operator apply，一个 logical task 负责一个 $r_i$，读取 $u$ 的邻域并只写自己的输出，因此
\begin{equation}
R_i=
\{u_j:j\in\mathcal N(i)\}
\cup
\{b_i,\text{operator data}\},
\qquad
W_i=\{r_i\}.
\end{equation}
这是标准的 unique-writer stencil gather。

如果整个所需邻域都在当前地址空间内，邻域读取只是普通 load；如果某些邻居由别的 rank 拥有，那么\textbf{同一个逻辑 task 并没有改变}，只是它的输入 readiness 需要 halo exchange 保证。可以把任务集合分为
\begin{equation}
\mathcal T
=
\mathcal T_{\mathrm{interior}}
\cup
\mathcal T_{\mathrm{boundary}},
\end{equation}
其中 interior task 的输入全部本地就绪，而 boundary task 还等待 ghost。这个分解是后面 MPI communication overlap 的算法来源，而不是 MPI API 自己创造出来的结构。

物理边界又是另一件事。边界 residual task 仍然产生同一个 $r_i$，只是算子行要使用\chapref{ch:discretization}推导的 Dirichlet、Neumann、Robin 或其他边界闭合。必须区分
\begin{equation}
\boxed{
\text{physical boundary condition}
\neq
\text{partition-boundary data readiness}.
}
\end{equation}

\section{Restriction 与 Prolongation 的关键是输出所有权}

\subsection{Restriction 作为 Coarse-Output Gather}

设 restriction 写成
\begin{equation}
b_H(I)
=
\sum_{j\in\mathcal S(I)}
w_{Ij}r_h(j).
\label{eq:mg-restriction-gather}
\end{equation}
令一个 coarse task 负责一个 $b_H(I)$，则输入是若干 fine residual，输出只有自己的 coarse cell。于是
\begin{equation}
W_I=\{b_H(I)\},
\end{equation}
天然满足 unique writer。

二维 full-weighting 一个 coarse output 最多读取 $3\times3$ 邻域，三维 tensor-product 形式最多读取 $3^3$ 个 fine values。这里维数提升影响的是\textbf{每 task 输入 footprint 与数据复用}，而不是 restriction 是否可并行。

如果 fine/coarse ownership 跨 rank 改变，restriction 还会附带 redistribution，但这是\textbf{输出所有权变化}引起的数据移动，不应与局部 restriction 数值公式混成一个概念。

\subsection{Prolongation 优先采用 Fine-Output Gather}

若按 coarse point 向多个 fine points scatter，
\begin{equation}
e_h(i)\mathrel{+}=p_{iI}e_H(I),
\end{equation}
多个 coarse tasks 可能同时贡献到同一个 fine output，从而形成 WAW 冲突。

更适合大规模并行的语义通常是让每个 fine task 主动 gather：
\begin{equation}
e_h(i)
=
\sum_{I\in\mathcal C(i)}
p_{iI}e_H(I).
\label{eq:mg-prolongation-gather}
\end{equation}
这样每个 fine output 仍然只有一个 writer。注意这里改变的是\textbf{计算组织方式}，不是数学上的 prolongation operator $P$。

如果 prolongation 后立即 correction，
\begin{equation}
u_h\leftarrow u_h+Pe_H,
\end{equation}
并且独立的 $e_h$ 不会再被后继操作读取，则可以把式~\eqref{eq:mg-prolongation-gather}直接融合进 correction，避免物化整个 fine correction 向量。这正是\chapref{ch:parallel-model}所说的“保持语义、改写 DAG 与存储”的例子。

\section{Correction、Reduction 与 Ghost Readiness}

Correction
\begin{equation}
u_h(i)\leftarrow u_h(i)+e_h(i)
\end{equation}
本身是纯 map；若 $e_h$ 已经就绪，每个输出可独立更新。若与 prolongation 融合，它则变成“coarse gather + fine map”的单 task。

残差范数
\begin{equation}
\|r\|_2^2=\sum_i r_i^2
\end{equation}
属于 reduction。其局部平方阶段具有 $\Theta(N)$ work 和大量并行度，真正的 span 来自聚合树。在共享内存中它映射为线程/树形归约，在 GPU 中变成 block/device reduction，在 MPI 中还需要 global collective。若 MG 只作为固定次数 preconditioner apply，内部不一定每层都需要 norm；\textbf{为了“监控一下”而增加的全局 norm 可能直接把一个本可局部执行的阶段变成全局同步点。}

Ghost fill 应理解为一个\textbf{readiness 操作}，而不是某种固定通信 API。它的后置条件是：
\begin{equation}
\boxed{
\text{依赖 ghost 的后继 stencil task 可以读取与 owner 一致的输入版本。}
}
\label{eq:mg-ghost-readiness}
\end{equation}
实现可能是：
\begin{itemize}
\item 同一进程内 patch-to-patch copy；
\item 周期边界映射；
\item AMR coarse-to-fine interpolation；
\item MPI halo exchange；
\item 以上几种操作的组合。
\end{itemize}
因此“ghost cell”是存储概念，“ghost readiness”才是 DAG 概念。

\section{V-Cycle 是一个嵌套 DAG}

把一次 V-cycle 简化成一条阶段链是有用的第一近似：
\begin{equation}
\boxed{
\begin{aligned}
&\text{pre-smooth}
\rightarrow
\text{residual}
\rightarrow
\text{restriction}
\rightarrow
\text{coarse cycle/solve}
\\
&\rightarrow
\text{prolongation/correction}
\rightarrow
\text{post-smooth}.
\end{aligned}
}
\label{eq:mg-parallel-vcycle-dag}
\end{equation}
但真正的执行图至少有三层结构：

\begin{enumerate}
\item \textbf{kernel 内部并行}：例如 residual 中大量独立 point/tile tasks；
\item \textbf{patch/box 级并行}：同一 level 上不同空间块可以独立推进到某些局部依赖边界；
\item \textbf{level 间递归关键路径}：fine residual 必须产生 coarse RHS，coarse error 又必须返回以后才能做 fine correction。
\end{enumerate}

第三层非常重要。典型 V-cycle 的 coarse solve 位于从 fine state 到最终 fine correction 的关键路径上，因此不能靠“fine level 有很多线程”隐藏一个完全串行的 coarse solve。多重网格的核心性能矛盾恰恰是：
\begin{equation}
\boxed{
\text{越粗的 level，数值 work 越少，}
\quad
\text{但它越容易成为同步与关键路径问题。}
}
\label{eq:mg-coarse-paradox}
\end{equation}

另一方面，式~\eqref{eq:mg-parallel-vcycle-dag}也不意味着每个阶段都必须有全局 barrier。若两个 patch 的依赖是局部的，某个 patch 的 residual/transfer 可以在它所需数据就绪后继续，而不必等待整个 level 所有 patch 完成。是否值得采用这种更细 DAG，要由 runtime 开销、实现复杂度与可重叠 work 决定。

\section{粗化怎样系统性耗尽并行度}

$d$ 维 full coarsening 下，未知量数近似满足
\begin{equation}
N_{\ell+1}
\approx
\frac{N_\ell}{2^d}.
\label{eq:mg-parallel-coarsening-count}
\end{equation}
若一个局部 kernel 每 unknown 的 work 为常数量级，
\begin{equation}
W_\ell=\Theta(N_\ell),
\end{equation}
那么各层总 work 构成几何级数。例如三维时
\begin{equation}
N_0+N_1+N_2+\cdots
\approx
N_0
\left(
1+\frac18+\frac1{8^2}+\cdots
\right)
=
\frac87 N_0.
\end{equation}
这解释了为什么从纯算术工作量看，增加很多 coarse levels 仍可能保持近线性复杂度。

但\textbf{work 小不等于时间一定小}。设调度时希望每个活跃 worker 至少获得 $g_{\min}$ 个 unknowns，则第 $\ell$ 层最多有
\begin{equation}
P_\ell^{\mathrm{useful}}
\lesssim
\min\left(
P,
\frac{N_\ell}{g_{\min}}
\right)
\label{eq:mg-useful-workers}
\end{equation}
个有意义的并行 worker。三维每粗化一次，$N_\ell$ 约缩小到原来的 $1/8$，可用 worker 数也可能近似按同样比例下降。

于是 coarse-level 时间更适合用
\begin{equation}
T_\ell
\approx
T_{\ell,\mathrm{useful}}
+
T_{\ell,\mathrm{fixed}}
+
T_{\ell,\mathrm{imbalance}}
\end{equation}
理解。$T_{\ell,\mathrm{useful}}$ 随 $N_\ell$ 快速下降，但 barrier、kernel launch、MPI latency、collective latency、调度和 metadata 等 $T_{\ell,\mathrm{fixed}}$ 不会同比例下降。最终会进入
\begin{equation}
\boxed{
T_{\ell,\mathrm{fixed}}
\gtrsim
T_{\ell,\mathrm{useful}}
}
\end{equation}
的区域，这就是 coarse-grid parallelism collapse。

后面的共享内存线程收缩、GPU coarse-kernel 合并、MPI rank consolidation、agglomeration 与 bottom solver 切换，本质上都是在处理式~\eqref{eq:mg-useful-workers}导致的同一个算法事实。

\section{数据生命周期决定 Fusion 与 Buffer Reuse}

MG 会产生多个临时状态，但“数学对象存在”不等于“整个 solve 期间都必须独占一块数组”。可以把每个向量看成具有
\begin{equation}
\text{birth}
\rightarrow
\text{last use}
\rightarrow
\text{storage reusable}
\end{equation}
的生命周期。

典型例子包括：
\begin{itemize}
\item fine residual $r_\ell$ 在 restriction 完成并确认没有后继读取以后可以复用；
\item prolongated correction 若只用于立即更新 $u_\ell$，可以不独立物化；
\item Jacobi 的 $u^{old}$ 与 $u^{new}$ 可以在 sweep 之间交换角色；
\item coarse temporary 在返回 fine level、最后一次 prolongation 完成后可以释放或进入池化复用。
\end{itemize}

Buffer reuse 的正确判定不是“这个数组现在看起来不用了”，而是
\begin{equation}
\boxed{
\text{所有读取该版本的后继 task 都已经完成。}
}
\end{equation}
过早覆盖会制造真正的 RAW/WAR 错误；过度保守则会增加峰值内存和 memory traffic。

Kernel fusion 也遵循同一规则。例如 residual 后立即 restriction 可以在某些局部算法中尝试 producer--consumer fusion，但只有在 coarse output ownership、fine residual 的其他用途和边界数据 readiness 都允许时才成立。\textbf{Fusion 是 DAG 重写，不是简单把两个 for-loop 拼在一起。}

\section{预条件 Krylov 中存在两层同步结构}

若固定 V-cycle 被用作近似逆
\begin{equation}
z=B_{MG}r,
\end{equation}
则它在外层 Krylov 方法中只是一个 preconditioner apply 子图：
\begin{equation}
\boxed{
\text{Krylov operator/dot}
\rightarrow
\underbrace{\text{MG V-cycle}}_{z=B_{MG}r}
\rightarrow
\text{Krylov update/dot}.
}
\label{eq:mg-parallel-preconditioned-krylov-dag}
\end{equation}

这时需要区分两类同步：

\begin{itemize}
\item MG 内部主要来自 stencil readiness、颜色阶段、transfer、coarse recursion 和可能的 coarse redistribution；
\item Krylov 外层还会引入 dot product、norm 等全局 reduction。
\end{itemize}

因此一个 MG kernel 即使局部扩展很好，整个 Krylov solve 仍可能被外层 global reductions 限制。反过来，为了减少 V-cycle 内部同步而改变 smoother，也必须重新检查 preconditioner quality，因为迭代次数可能上升。

如果 $B_{MG}$ 每次 apply 都执行固定 cycle、固定 smoother 和固定粗层策略，可以把它看成固定线性算子。若 inner tolerance、cycle 数或粗层策略随迭代变化，则 preconditioner apply 本身可能变化，需要\secref{sec:mg-preconditioner}中的 flexible Krylov 视角。

\section{维数、边界与分区改变的是成本结构}

对边长约为 $n$ 的 $d$ 维规则子域，内部 work 与体积成正比，
\begin{equation}
W_{\mathrm{interior}}\sim n^d,
\end{equation}
而一层边界/halo 数据与表面积成正比，
\begin{equation}
Q_{\mathrm{boundary}}\sim n^{d-1}.
\end{equation}
因此
\begin{equation}
\frac{Q_{\mathrm{boundary}}}{W_{\mathrm{interior}}}
\sim
\frac1n.
\label{eq:mg-parallel-surface-volume}
\end{equation}
这一个几何关系会在不同硬件上表现成不同成本：
\begin{itemize}
\item CPU tile 边界主要造成 cache/workset 重复与可能的 NUMA 访问；
\item GPU tile/block 边界影响 shared-memory reuse、global-memory traffic 与 occupancy；
\item MPI partition boundary 才真正产生远端消息。
\end{itemize}

还必须区分三种“边界”：
\begin{equation}
\boxed{
\text{physical boundary}
\neq
\text{partition boundary}
\neq
\text{AMR coarse--fine interface}.
}
\label{eq:mg-boundary-classes}
\end{equation}
物理边界改变离散算子；partition boundary 不改变 PDE，只改变输入数据的取得方式；AMR coarse--fine interface 则需要层间插值、平均或通量匹配。三者都可能使用 ghost-like storage，却来自完全不同的数学语义。

\section{MG 操作的统一并行画像}

\begin{center}
\small
\begin{tabularx}{0.99\textwidth}{>{\bfseries}p{2.0cm}p{2.35cm}p{2.65cm}X}
\toprule
操作 & 典型数据流 & 核心状态依赖 & 主要并行风险/成本来源\\
\midrule
Jacobi & stencil gather + map & $u^{(k)}\rightarrow u^{(k+1)}$ & old/new 存储、带宽、sweep boundary\\
GS/RBGS & ordered/colored stencil & 同 sweep 新值被后继读取 & 长依赖链、color barrier、颜色数\\
算子/残差 & stencil gather & 输入版本/ghost ready & bandwidth、boundary readiness\\
restriction & coarse-output gather & fine residual ready & input footprint、ownership redistribution\\
prolongation & fine-output gather & coarse error ready & scatter 冲突、coarse/fine ownership\\
correction & map 或 fused gather & correction ready & memory pass、fusion lifetime\\
norm/dot & reduction & 全部局部贡献 ready & reduction span、global collective\\
ghost fill & copy/interpolate/exchange & owner version $\rightarrow$ ghost version & copy、插值、message latency\\
coarse solve & 依算法而定 & coarse RHS ready & 并行度塌缩、固定开销、bottom solver\\
\bottomrule
\end{tabularx}
\end{center}

这张表不是实现菜单，而是后续硬件章节的共同语义接口。任何 OpenMP、CUDA、MPI 或混合实现都应该能说明自己如何实现这些状态依赖和所有权规则。

\section{从一个新 MG Kernel 到可执行设计}

遇到新的 MG 操作或优化方案，可以按以下顺序判断：

\begin{enumerate}
\item 先写数学输入、输出和状态版本，不看现有循环代码；
\item 选择 point/line/tile/patch 级 logical task，列出 $R_t,W_t$；
\item 判断依赖是 RAW 真依赖，还是由 in-place storage 引入的 WAR/WAW；
\item 尽量寻找 unique-writer gather 表达；
\item 确定哪些 task 只依赖 local/owned 数据，哪些要等 ghost、coarse state 或 reduction；
\item 估算本层 $W_\ell$、可用 task 数和合理的 $P_\ell^{\mathrm{useful}}$；
\item 再估 bytes、temporary lifetime、message volume、barrier/reduction count；
\item 最后才选择线程、GPU 或 rank 映射，并用 profiler 验证瓶颈。
\end{enumerate}

如果某个优化无法在这条链上说明“它究竟改变了 DAG、数据移动、粒度还是固定开销”，那么它通常还没有被充分理解。

\section{本章小结}

多重网格的并行结构不是“把每个 for-loop 加上 parallel”这么简单。Jacobi 的并行性来自 old/new 状态分离；Gauss--Seidel 的障碍来自新值立即可见；restriction/prolongation 的关键是输出所有权；norm/dot 的关键是聚合 span；ghost fill 的本质是建立输入版本 readiness；完整 V-cycle 则是一个带 level recursion 的嵌套 DAG。

粗化同时带来两个方向相反的效果：总算术 work 形成快速收敛的几何级数，但可用 task 数也按 $2^d$ 的因子迅速减少。于是越到粗层，越不是“算得多”，而是“可并行工作不足以摊薄同步、通信与调度固定成本”。后续 CPU、GPU 与 MPI 章节只是把这一统一结构映射到不同机器。

\section*{习题}

\subsection*{状态与依赖推导}
\begin{enumerate}[label=\arabic*.]
\item 对 weighted Jacobi 写出 $u^{(k)}$、$r^{(k)}$、$u^{(k+1)}$ 三种状态之间的 DAG。分别画出“先 residual 后 update”和“fused Jacobi”两种组织，并指出后者消除了哪些中间数据流。
\item 对原地 lexicographic GS 的相邻三个点写出 RAW 依赖；再对标准二维五点 RBGS 画出 red/black 两阶段 DAG。
\item 说明为什么“二维五点 RBGS 可以两色”不能直接推广为“任何 GS 都只需要两色”，并给出一个含对角耦合的 stencil 作为反例。
\item 对 gather restriction 与 scatter prolongation 分别列出 $R_t,W_t$，指出何时会出现 WAW。
\item 把 ghost fill 表述为 producer--consumer 依赖，区分 owner version、ghost version 和后继 boundary stencil task。
\end{enumerate}

\subsection*{Work、并行度与粗层分析}
\begin{enumerate}[label=\arabic*.]
\item 对 $256^3\rightarrow128^3\rightarrow64^3\rightarrow32^3\rightarrow16^3$ 层级，计算每层 unknown 数与相对 finest level 的比例。
\item 证明三维 full coarsening 下 $\sum_\ell N_\ell\le(8/7)N_0$，并解释为什么这个结果不能推出 coarse levels 的 wall time 可以忽略。
\item 设 $P=64$，希望每个活跃 worker 至少处理 $g_{\min}=4096$ 个 unknowns。对上一题各层计算式~\eqref{eq:mg-useful-workers}给出的最大有用 worker 数。
\item 假设每层 useful work time 与 $N_\ell/P_\ell$ 成正比，而每层固定同步成本为 $5\,\mu s$。构造一个简单模型，找出从哪一层开始 fixed cost 超过 useful work。
\item 推导 $d$ 维 full coarsening 下总 work 的几何级数上界，并比较 $d=2$ 与 $d=3$ 的并行度衰减速度。
\end{enumerate}

\subsection*{实现与验证}
\begin{enumerate}[label=\arabic*.]
\item 实现 residual、Jacobi、restriction、prolongation、correction 和 norm 的串行版本，记录每个 kernel 的 FLOPs、读写字节和输出所有权。
\item 实现独立 $e_h$ 数组的 prolongation+correction 与 fused 版本，验证数值一致性并比较 memory traffic。
\item 为三层 V-cycle 记录每个临时向量的 birth、last use 与可复用时刻，设计一套不破坏依赖的 buffer pool。
\item 把 residual task 分成 interior 与 ghost-dependent boundary 两组，在单进程环境中先验证两组任务任意交错执行仍得到相同结果。
\item 对 fixed V-cycle 与“每层都计算 residual norm”的版本比较同步点数量，解释哪些 norm 是算法必需、哪些只是诊断开销。
\end{enumerate}

\subsection*{工程诊断}
\begin{enumerate}[label=\arabic*.]
\item 某 coarse level 只有 12000 unknowns，却仍启动 64 个 CPU threads。按本章模型列出可能的效率损失，并设计 thread shrinking 规则。
\item 某 GPU V-cycle 的 finest residual 已达到很高带宽，但总 solve 仍很慢。只使用本章抽象，列出 coarse-level parallelism、reduction、launch 与 data-lifetime 四条诊断路径。
\item 某 MPI 版本 restriction 后立刻做 rank consolidation。说明需要区分哪些时间属于数值 restriction，哪些属于 ownership redistribution，并给出 break-even 判断所需的测量项。
\item 一个新的 smoother 每 sweep 比 Jacobi 慢 2 倍，但 V-cycle 数减少到原来的 35\%。说明为什么只能用 time-to-solution 而不能用单 sweep 时间判断优劣，并列出还应比较的同步与并行度指标。
\end{enumerate}
